\homework{2}{Wed 01 Oct 2025 21:25}{}

\textbf{Free Particle in Minkowski Spacetime.} Consider a free particle $x^\mu(t)$ in the Minkowski spacetime $ds^2 = \eta_{\mu \nu} dx^\mu dx^\nu$. Take as its action $S = - m \int ds = - m \int dt \left( \frac{ds}{dt} \right) $. Determine its equation of motion.

\begin{proof}[Solution]
	We observe that $S$ only depends on $\dot{x}^\mu$, thus
	\begin{align*}
		\delta S
		&= - m \int dt \frac{\partial \sqrt{\eta_{\mu \nu}\dot{x}^\mu \dot{x} ^\nu} }{\partial \dot{x}^\alpha} \delta \dot{x}^\alpha \\
		&= - m \int dt \frac{\eta_{\alpha \beta}\dot{x}^\beta \delta\dot{x} ^\alpha} {2 \frac{ds}{dt}} \\
		&= 0
	.\end{align*}
	For this integral to be zero, we must have
	\[
		\eta_{\alpha \beta} \dot{x}^\beta \delta\dot{x} ^\alpha = 0
	.\] 
	Integration by part yields
	\[
		\eta_{\alpha \beta} \ddot{x} ^\beta \delta x^\alpha = 0
	.\] 
	Which implies
	\[
		\ddot x = 0
	.\] 
	This means that the particle undergoes unaccelerated motion.
\end{proof}


\textbf{Coulomb gauge}. Consider the classical Maxwell fields.
\begin{itemize}
	\item Expand the Lagrangian density in the presence of a current in terms of $A^\mu = (\varphi, \vec{A} )$ with $\varphi$ the electrostatic potential and $\vec{A} $ satisfying the coulomb gauge $\vec{\nabla} \cdot \vec{A} (x) = 0$.
	\item Vary the action with respect to $\varphi$ and $\vec{A} $ to find their equations of motion.
	\item Find the solution to the equations of motion using Fourier transform.
\end{itemize}
\begin{proof}[Solution]
	Recall the least action principle for classical field theory:
	\[
		\delta S = \delta \int \mathcal{L} d x^\mu = 0
	.\] 
	where we have for electromagnetism,
	\[
		\mathcal{L}(\partial_\mu A_\nu, A_\nu)
		= - \frac{1}{4} F_{\mu \nu} F^{\mu \nu} + J_\nu A^\nu
	\] 
	where $F_{\mu \nu} = \partial_\mu A_\nu - \partial_\nu A_\mu$,  $J_\nu = (\rho, \vec{j} )$, and $A_\nu = (\varphi, \vec{A} )$.

	Vary the Lagrangian, we find
	\begin{align*}
		\delta S
		&= \int d^4 x - F^{\mu \nu} \delta \left( \partial _\mu A_\nu \right) + J_\nu \delta A^\nu \\
		&= \int d^4 x \partial _\mu F^{\mu \nu} \delta A_\nu  + J_\nu \delta A^\nu \\
		&= \int d^4 x \left(\partial _\mu F^{\mu \nu}  + J_\nu\right) \delta A^\nu = 0
	.\end{align*}
	Thus, the field-theoretic equation of motion is
	\[
		\partial _\mu F^{\mu \nu} = J^\nu
	.\] 

	Under the Coulomb gauge $\nabla \cdot  \vec{A} = 0$, the equation becomes
	\begin{align*}
		\Delta \varphi &= - \rho\\
		\square \vec{A}  &= - \vec{J} + \nabla  \partial_t \varphi
	.\end{align*}
	The first equation is Poisson equation for electrostatics. It is solved via the Green's function in three dimensions
	\[
		G_3(x) = \frac{1}{4 \pi |x|}
	.\] 
	Thus the solution is
	\[
		\varphi(t, x) = \int d^3 x' \frac{- \rho(t, x')}{4 \pi |x - x'|}
	.\] 
	The second equation is three decoupled wave equations with source terms. Solution is given by the Green's function
	\[
		G(t, x) = \frac{1}{4 \pi |x|} \delta\left(t-|x|\right) 
	.\] 
	This is derived by Fourier-transforming the wave equation and choosing an appropriate integration contour that preserves causality.
	The solution of the wave part reads
	\[
		\vec{A} (t, x) = 
		int dt' \int d^3 x \frac{1}{4 \pi |x - x'|} \delta\left(t-t'-|x-x'|\right) \left( \nabla  \partial_t \varphi - \vec{J}  \right) 
	.\] 

\end{proof}


\textbf{Spontaneous Symmetry Breaking.} Consider a complex relativistic scalar field with Lagrangian density
\[
	\mathcal{L} = - \frac{1}{2} \left( \partial_\mu \Phi^\dag \right) \left( \partial^\mu \Phi \right) 
	+ \lambda \left( \Phi^\dag \Phi - \phi_0 \right) ^2
.\] 
\begin{itemize}
	\item Find the classical equations of motion
	\item Consider constant field configurations and find minima of action.
	\item Expand the field $\Phi(x^\mu)$ around this minimum to the second order and find the mass of the excitations.
	\item Is there a massless mode?
\end{itemize}

\begin{proof}[Solution]
	Start by computing the Euler-Lagrange equation
	\[
		\delta S = 0 \implies
		\frac{\partial \mathcal{L}}{\partial \Phi}  - \partial_\mu \frac{\partial \mathcal{L}}{\partial (\partial_\mu \Phi)}  = 0
	.\] 
	\begin{align*}
		\frac{\partial \mathcal{L}}{\partial \Phi} &= 2 \lambda \Phi^\dag \left( \Phi^\dag \Phi - \phi_0 \right) \\
		\frac{\partial \mathcal{L}}{\partial (\partial_\mu \Phi)} &= - \frac{1}{2} (\partial^\mu \Phi)^\dag = - \frac{1}{2} \partial^\mu \Phi^\dag \\
		\partial_\mu \frac{\partial \mathcal{L}}{\partial (\partial_\mu \Phi)} &= - \frac{1}{2} \partial_\mu \partial^\mu \Phi^\dag
	.\end{align*}
	The equation of motion reads
	\[
		( \square + 4 \lambda \left( \Phi^\dag \Phi - \phi_0 \right) ) \Phi = 0
	.\] 
	This looks like the Klein-Gordon equation except with an extra mass term $\Phi^\dag \Phi$.

	Under the assumption of constant field, $\partial _\mu \Phi \equiv 0$, the Lagrangian density simplifies to
	\[
		\mathcal{L} = \lambda \left( \Phi^\dag \Phi - \phi_0 \right) ^2
	.\] 
	Assume $\lambda > 0, \phi_0 \ge 0$. We see that the Lagrangian density acquires minimal value $0$ when $\Phi^\dag \Phi = \phi_0$. Hence, actions achieves minimum value $S = 0$ when $\Phi = \sqrt{\phi_0} e^{i \alpha}$, for some $\alpha > 0$.

	To expand $\Phi$ around this minimum, put
	\[
		\Phi(x) = \sqrt{\phi_0}  e^{i \alpha} + h(x) + i G(x)
	\] 
	where $h, G$ are real scalar fields close to $0$. Insert in the full Lagrangian density and keep only up to second order, we see that
	\[
		\mathcal{L} = - \frac{1}{2}( \partial_\mu \partial^\mu h + \partial_\mu \partial^\mu G) 
		+ 4 \lambda \phi_0 \cos^2 \alpha h^2 + h.c.
	.\] 
	There is no mass term for field  $G$; this implies existence of a massless mode (Goldstone boson) described by the field $G(x)$, which is decoupled (up to second order) from $h$.
\end{proof}


